\chapter{Como ler este livro}
\label{cap:como-ler}

\rascunho{
  Mapa do livro e roteiros de leitura por perfil. Ideias:
  \begin{itemize}
    \item As cinco partes em um parágrafo cada (\nameref{parte:origens}, \nameref{parte:trilhas}, \nameref{parte:modelagem}, \nameref{parte:devops}, \nameref{parte:agil}).
    \item Roteiros: quem está começando lê tudo em ordem; quem já trabalha com dados pode ir direto à Parte~\ref{parte:modelagem}, \ref{parte:devops} ou \ref{parte:agil}; gestores podem ler a Parte~\ref{parte:origens} e a Parte~\ref{parte:agil}.
    \item Convenções do livro: as caixas \enquote{Da minha experiência} são relatos de projetos reais (sem nomes de empresas nem clientes); \enquote{Antes e depois da IA} compara o mesmo trabalho nos dois momentos; trechos de código em SQL e Python.
  \end{itemize}
}

\pendente{escrever o guia de leitura}
